\section{Introduction}
\label{sec:intro}
A person using a symbol-based augmentative and alternative communication (AAC) board selects pictograms to express a request, refusal, feeling, or action.
The selected labels are often telegraphic: \emph{Water}, \emph{No + Pizza}, or \emph{Tired}.
Expanding these labels requires grammatical material, but an added meal, changed speaker, or missing negation can misrepresent the person whose message is being voiced.
The central problem is to permit sentence realization while making the acceptance decision inspectable.

\begin{figure}[t]
\centering
\setlength{\fboxsep}{6pt}
\fcolorbox{exampleblue}{panelgray}{\begin{minipage}{0.91\columnwidth}\small
\textbf{1\quad Select}\enspace Water\quad[request]\par\medskip
\textbf{2\quad Propose}\enspace I want water \bad{today}.\par
\hfill\bad{[added: time]}\par\medskip
\textbf{3\quad Check}\enspace concept $+$ polarity $+$ act\par
\hfill reject the added time\par\medskip
\textbf{4\quad Return}\enspace\good{I want water.}\par
\hfill checked template\par\medskip
\footnotesize If no sentence passes: show literal words.\par
If the act is ambiguous: ask the user.
\end{minipage}}

\caption{Concept-gated realization (authored illustration). A water request is expanded, checked, and either accepted or replaced by a checked template. The red, labelled addition is unselected content; blue marks the retained message. A literal preview is shown when neither path passes. These are decision rules, not a semantic guarantee.}
\label{fig:teaser}
\end{figure}

For a water request, a single symbol can be realized as a complete sentence in each language.
The following are authored illustrations, with \lang{English} as a reference language and the three Indic languages as the proposed extensions:
\begin{quote}\small
\lang{English}\enspace\good{I want water.}\\
\lang{Hindi}\enspace\good{\hi{मुझे पानी चाहिए।}}\\
\lang{Gujarati}\enspace\good{\gu{મને પાણી જોઈએ છે.}}\\
\lang{Marathi}\enspace\good{\hi{मला पाणी हवे आहे.}}
\end{quote}
\noindent English gloss: \emph{“I want water.”}
The language-specific predicates and negation patterns motivate checks beyond an English-only word list.
A single symbol does not determine its act in every context; the water example assumes a request, rather than inferring that intention universally.

\sys{} combines a symbol board, a prompted \emph{\model{}} model \citep{qwen2024model}, a \emph{concept ledger}, deterministic acceptance checks, and language-specific fallback templates (Figure~\ref{fig:teaser}).
The ledger records selected concepts, the intended act, and polarity; the checker tests whether a candidate conforms to those constraints.
The revised interface separates act from negation and asks for clarification when the selection has no determinate act.
The intended audience is symbol-board users and the caregivers, clinicians, and researchers who configure their vocabulary.

Our novelty is the design and audit of explicit acceptance checks in this multilingual workflow, including what the fallback policy does when generation fails.
Grammatical expansion and Indian-language AAC already exist (Section~\ref{sec:related}); we do not claim either as a first.
The archived web-client comparison improves lexical success by 18.3 percentage points over its ungated cascade, but remains 12.0 points below the template baseline.
This trade-off motivates measuring model coverage and failures alongside success, rather than treating every accepted sentence as faithful.
We make three scoped contributions:
\begin{itemize}
\item \textbf{An inspectable realization policy} connecting the concept ledger, candidate checks, retry, template fallback, clarification, and literal preview.
\item \textbf{A paired four-language audit} that separates browser interpretation, model output, checker decisions, templates, and scorer defects.
\item \textbf{Evidence-linked revisions and examples} that document the repairs, an explicitly labelled regression comparison, and representative input--output cases in all four languages.
\end{itemize}
